\section{公共求解器、验证协议与实验设计}
\subsection{三层评估结构}
公共求解器将每个候选方案依次送入三个层次。轻量层使用驻留域字节、依赖下界、Pipe 工作量和同步等待构造可复算的排序分数；展开层调用题面核内步骤，得到确定性拓扑序、容量峰值和 Spill 记录；完整层调用官方多核事件模拟，得到 Makespan、额外搬运和 Cache 字段。轻量层和展开层的时间只能说明候选排序或核内结构，不能代替完整结果。

\subsection{规范化与去重}
规范化只重新编号子图，不改变每个核心列表中的相对顺序。计划键由图 id、场景、核数、配置哈希、官方代码哈希和规范化方案共同构成。相同计划在相同输入下只做一次对应层级评估；失败计划也保存失败类别，避免重复消耗并掩盖失败分布。

\subsection{离散事件模拟的状态}
每个操作依次处于 `pending`、`ready`、`running` 和 `done` 状态。进入 `ready` 需要满足所属 Task 已激活、核内前驱完成、跨核释放时刻已到，并在需要申请核内存储时满足申请序和容量条件。发射时更新 Pipe、张量引用、共享 DDR/Cache 带宽池和预计完成时刻；时间推进取下一完成事件与下一跨核释放事件的最小值。该过程直接对应题面评估器的状态语义，不在论文中把核内参考时间写成最终 Makespan。

\begin{algorithm}[H]\caption{公共候选评估流程}\begin{algorithmic}[1]
\REQUIRE 图 $G$、场景 $r$、核数 $K$、配置和候选生成规则
\ENSURE 候选账本、官方结果或失败状态
\STATE 解析图并建立张量消费者、依赖和驻留域索引
\STATE 生成并规范化起点，删除重复计划
\STATE 计算轻量分数并保存候选来源
\STATE 对有限候选执行核内展开，记录容量、顺序和失败原因
\STATE 对入选计划执行官方多核事件模拟
\STATE 保存 Makespan、搬运、Cache 字段、计划哈希和运行状态
\STATE 生成可供论文读取的表格和图表输入，但不在运行中写入正文
\end{algorithmic}\end{algorithm}

\subsection{实验矩阵与结果门槛}
主矩阵覆盖题面规定的全部图、1 至 5 核和 A/B/C 场景；单核基线由统一核内调度在整图上计算。实验脚本、运行号、参数、环境和代码哈希由 `华为杯\_code` 保存，冻结工作区通过符号链接只读查看。当前主矩阵仍在运行，本稿不读取其数值生成结论。

实验结果进入正文前必须同时满足：运行状态为可解释的终态；结果字段通过程序检查；主矩阵与单核基线的样本和指标口径一致；逐用例缺失、失败和 PROFILE 状态有明确分母；图表由结果 CSV 重建；队伍完成科学复核。任何未满足的项在论文中保持 `\pending` 或“待补当前运行结果”。

\subsection{稳健性和失败分析计划}
正式结果完成后，至少检查以下内容：小规模可解析图对拍；不同核数下的方案合法性和约束回代；容量边界与退化到单核/无 Cache 的行为；候选失败与超时分布；同一计划重复读取的一致性；Cache 命中率、DDR 服务字节和 Makespan 的方向是否相互可解释。随机性只在代码确实使用随机过程时报告重复波动；当前确定性候选排序不虚构置信区间。
